\FloatBarrier
\section{Calculating Many-Electron Energies}
\label{app:interacting-models}
\label{sec:interacting-electron-models}

Section~\ref{sec:quantum-dot-introduction} defines $H_N$ and its ground
energy $E_N$; Section~\ref{sec:electron-number-control} compares these
energies to determine the electron number. This appendix calculates
the many-electron energies in an orbital basis, using the screened
pair interaction $W$ from Appendix~\ref{app:screened-interaction}.
The localized basis and its matrix elements also give the several-dot description in
Section~\ref{sec:multiple-quantum-dots}, including the tunnel matrix
in Eq.~\eqref{eq:dot-pair-hamiltonian}.
The orbitals are those described in
Appendix~\ref{app:effective-mass}, with valley components from
Appendix~\ref{app:valley-states} when they need to be retained.

With two electrons in the device, their mutual repulsion adds
an energy that depends on both positions. For $N$ electrons,
we therefore use
\begin{equation}
 H_N=\sum_{a=1}^{N}h_1(a)
       +\sum_{a<b}W(\mathbf r_a,\mathbf r_b).
 \label{eq:device-continuum-many-electron}
\end{equation}
Here $h_1(a)$ acts on electron $a$ and includes kinetic energy,
external confinement, and any retained valley and spin terms.
$W$ is the interaction energy of a pair. The restriction $a<b$
counts each pair once and excludes interaction of an electron
with itself. $H_N$ acts on antisymmetric $N$-electron states.

\subsection{Self-consistent mean-field approximation: Schr\"odinger--Poisson}
\label{sec:device-schrodinger-poisson}

One way to approximate the interaction term in
Eq.~\eqref{eq:device-continuum-many-electron} is to replace the
electrons' position-dependent repulsion by the electrostatic
potential of their average charge density. If $n_q(\mathbf r)$
is the electron number density, this gives the Hartree potential
energy and the mean-field one-particle operator
\begin{equation}
 V_{\mathrm H}[n_q](\mathbf r)
       =\int W(\mathbf r,\mathbf r')n_q(\mathbf r')\,d^3r',
 \qquad
 h_{\mathrm{MF}}[n_q]=h_1+V_{\mathrm H}[n_q].
 \label{eq:device-hartree-mean-field}
\end{equation}
The density determines the potential, while the occupied states
in that potential determine the density. A self-consistent solution
must satisfy both requirements simultaneously.

Assume spin-degenerate levels. Let $\mathbf F_a$ be a normalized
multivalley orbital with energy $\epsilon_a^{\mathrm{MF}}$, and
let $p_a$ be its mean occupation, counting both spins. At
temperature $T$, with chemical potential $\mu$ and Boltzmann
constant $k_B$, the \emph{Schr\"odinger--Poisson equations} are
\begin{equation}
 \left\{
 \begin{aligned}
  [h_1+V_{\mathrm H}[n_q]]\mathbf F_a
       &=\epsilon_a^{\mathrm{MF}}\mathbf F_a,\\[3pt]
  n_q(\mathbf r)
       &=\sum_a p_a\sum_\nu|F_{a\nu}(\mathbf r)|^2,\\[3pt]
  p_a&=\frac{2}{1+\exp[(\epsilon_a^{\mathrm{MF}}-\mu)/(k_BT)]}.
 \end{aligned}
 \right.
 \label{eq:device-sp-system}
\end{equation}
The unknowns are the orbitals, orbital energies, density, and
occupations. Choose the orbitals orthonormal, with
$\sum_\nu\int|F_{a\nu}|^2\,d^3r=1$, and use the envelope
boundary and interface conditions of $h_1$.
Set $n_q=0$ outside the region used for the envelope calculation.

The occupation condition closes the problem in one of two ways.
For a prescribed mean electron number $N$, determine $\mu$
along with the other unknowns by imposing
\begin{equation}
 \sum_a p_a=\int n_q(\mathbf r)\,d^3r=N.
 \label{eq:device-sp-number}
\end{equation}
At zero temperature, fill the lowest self-consistent levels
subject to this number constraint. For equilibrium with an
electron reservoir, instead set $\mu=\mu_{\mathrm R}$; then
$\sum_a p_a$ is an output. Spin-resolved levels use individual
occupations with numerator one in the last line of
Eq.~\eqref{eq:device-sp-system}.

To solve these equations, start with a trial density, calculate
$V_{\mathrm H}[n_q]$ by integration using
Eq.~\eqref{eq:device-hartree-mean-field},
find the orbitals of $h_1+V_{\mathrm H}[n_q]$, and update their
occupations and density. At convergence, the density reconstructed
from the orbitals is the same density used to calculate the
Hartree potential~\cite{QTCADSchrodingerPoisson}.

This electrostatic mean-field treatment retains the Hartree
potential and omits exchange and correlation.\footnote{Using the
total density in $V_{\mathrm H}$ also introduces a spurious
self-interaction unless it is explicitly corrected. In particular,
the pair-interaction term of $H_N$ vanishes for $N=1$.}

After convergence, a finite set of occupied and unoccupied
orbitals can also be used to construct an antisymmetric basis
for a subsequent many-electron calculation. Denote the
electron number used to generate these orbitals by $N_{\mathrm{ref}}$.
For a target $N$-electron problem, one may choose
$N_{\mathrm{ref}}=N$; to compare several electron numbers, a common
reference basis can be used for all of them. The reference number
selects the orbital shapes, not the electron number of the later
calculation. Finite-basis results can depend on this choice, so
check convergence as more orbitals are retained.


\subsection{Basis choice and second quantization}
\label{sec:device-orbital-basis}

To write $H_N$ in second quantization, choose a finite set of
orthonormal one-electron spatial states $w_i(\mathbf r)$,
$i=1,\ldots,M$. Several states can be retained on each dot.
We neglect spin mixing. A spatial state $w_i$ may already include
local valley coupling, as in Section~\ref{sec:multiple-quantum-dots};
its matrix elements then use the full orbital--valley inner product.
In the single-valley limit, $w_i$ is a normalized scalar envelope
with the common crystal Bloch factor suppressed. The scalar integral
expressions below refer to that limit.
The states $w_i$ may be obtained directly from $h_1$ or from
the self-consistent calculation above. A dot model is easier to
interpret if each state is concentrated near a particular dot.
When using states from separate isolated-well calculations,
orthonormalize them first.

Consider two neighboring wells in a fixed spin sector. Their exact one-electron
eigenstates generally spread over both wells. Linear combinations
within their low-energy subspace can instead give states $w_1,w_2$
localized mainly on the left and right. In that basis,
\begin{equation}
 h_{ij}=\bra{w_i}h_1\ket{w_j},
 \qquad
 h_{\mathrm{pair}}=
 \begin{pmatrix}
  \epsilon_1&-t_{12}\\
  -t_{12}^*&\epsilon_2
 \end{pmatrix},
 \quad
 \epsilon_i=h_{ii},\quad t_{12}=-h_{12}.
 \label{eq:device-one-body-matrix}
\end{equation}
For scalar envelopes,
$h_{ij}=\int d^3r\,w_i^*(\mathbf r)h_1w_j(\mathbf r)$;
if valley components are retained, the inner product also sums
over those components.
The diagonal entries are the energies of the chosen localized
states. The off-diagonal entry coherently transfers amplitude
between them: this is the tunnel coupling.

Diagonalizing the matrix gives the level separation
\begin{equation}
 \Delta E=\sqrt{(\epsilon_1-\epsilon_2)^2+4|t_{12}|^2}.
 \label{eq:device-two-level-splitting}
\end{equation}
At equal dot energies, $\Delta E=2|t_{12}|$.
Away from this point, both the energy offset and tunnel coupling
contribute to the splitting.

For self-consistent orbitals, compute the same matrix elements
from $h_1$. The interaction between the dot electrons will be
included explicitly through $W$, so their Hartree potential is
not part of the one-body matrix.

In the same basis, the pair matrix element is
$\mathcal V_{ijkl}=\langle w_iw_j|W|w_kw_l\rangle$, with products
referring to the two electron coordinates. For scalar envelopes,
\begin{equation}
 \mathcal V_{ijkl}
   =\int d^3r\,d^3r'\,
      w_i^*(\mathbf r)w_j^*(\mathbf r')W(\mathbf r,\mathbf r')
      w_k(\mathbf r)w_l(\mathbf r').
 \label{eq:device-coulomb-matrix}
\end{equation}
This coefficient connects a pair initially in spatial states
$k,l$ to a pair in states $i,j$. Different choices of these indices
describe charge repulsion or transitions between orbital configurations.

With these one-body and two-body matrix elements, we can write
the many-electron Hamiltonian using
$c_{i\sigma}^\dagger$ and $c_{i\sigma}$, which add or remove an
electron in spatial state $i$ with spin
$\sigma\in\{\uparrow,\downarrow\}$. Their fermionic
anticommutation rules enforce the required antisymmetry.
Retaining all matrix elements gives
\begin{equation}
 H^{(M)}
 =\sum_{ij,\sigma}h_{ij}c_{i\sigma}^\dagger c_{j\sigma}
  +\frac12\sum_{ijkl,\sigma\sigma'}
      \mathcal V_{ijkl}
      c_{i\sigma}^\dagger c_{j\sigma'}^\dagger
      c_{l\sigma'}c_{k\sigma}.
 \label{eq:device-projected-hamiltonian}
\end{equation}
The factor $1/2$ avoids counting the same electron pair twice.
This is the full one-body-plus-pair Hamiltonian within the selected
orbital space~\cite{QTCADManyBody}. If spin mixing is included,
the one-body term becomes
$\sum_{ij,\sigma\sigma'}h_{ij}^{\sigma\sigma'}
c_{i\sigma}^\dagger c_{j\sigma'}$. Additional valley states can
likewise be included in the state index.

Within the retained orbital space, a unitary basis change
transforms both $h_{ij}$ and $\mathcal V_{ijkl}$ while preserving
the spectrum of $H^{(M)}$. Individual matrix elements refer to a
specified basis.

\paragraph{Grouping states by dot occupation.}
In a localized basis, assign each spatial state to a dot. If
$\mathcal I_d$ is the set of orbital indices assigned to dot $d$,
its number operator is
\[
 \hat N_d=\sum_{i\in\mathcal I_d,\sigma}
                 c_{i\sigma}^{\dagger}c_{i\sigma}.
\]
At fixed total $N$, let $P_{\mathbf N}$ project onto the common
eigenspace of these operators with occupations
$\mathbf N=(N_1,N_2,\ldots)$, where $N_d$ is the eigenvalue of
$\hat N_d$ and $\sum_d N_d=N$. These projectors satisfy
\begin{equation}
 P_{\mathbf N}P_{\mathbf N'}
   =\delta_{\mathbf N\mathbf N'}P_{\mathbf N},
 \qquad
 \sum_{\mathbf N}P_{\mathbf N}=I_N,
 \label{eq:device-charge-projectors}
\end{equation}
where the sum includes all allowed occupation tuples at the fixed
total $N$, and $I_N$ is the identity in the retained $N$-electron space.
Inserting this resolution of the identity on both sides of
$H^{(M)}$ gives the block decomposition used in
Eq.~\eqref{eq:dot-charge-blocks}:
\begin{equation}
 \begin{aligned}
 H^{(M)}&=H^{(0)}+\mathcal T,\\
 H^{(0)}&=\sum_{\mathbf N}P_{\mathbf N}H^{(M)}P_{\mathbf N},
 &\mathcal T&=\sum_{\mathbf N\ne\mathbf N'}
            P_{\mathbf N}H^{(M)}P_{\mathbf N'}.
 \end{aligned}
 \label{eq:device-charge-blocks}
\end{equation}
The block $H_{\mathbf N}$, written explicitly as
$H_{(N_1,N_2,\ldots)}$ in Eq.~\eqref{eq:dot-charge-blocks}, is the restriction of
$P_{\mathbf N}H^{(M)}P_{\mathbf N}$ to its charge subspace.
Diagonalizing each block separately gives its ground and excited
energies $E_{\mathbf N,a}$. All terms that preserve every dot
occupation remain in these blocks, including interdot repulsion
and direct Coulomb exchange. The off-diagonal operator $\mathcal T$
includes single-electron hopping and any interaction matrix elements
that change dot occupations, such as pair hopping. Dropping
$\mathcal T$ therefore defines a reference with no charge transfer;
it still retains interactions between different dots.

If $\ket{\mathbf N,a}$ and $\ket{\mathbf N',b}$ are nondegenerate
block eigenstates, first-order perturbation theory gives the
admixture of the latter into the former as
\[
 \frac{\bra{\mathbf N',b}\mathcal T\ket{\mathbf N,a}}
      {E_{\mathbf N,a}-E_{\mathbf N',b}}.
\]
This expression explains why the relevant comparison is with the
energy separation of the connected states. When that separation
is comparable to their coupling, retain both states and diagonalize
the coupled block. For the one-electron two-dot matrix in
Eq.~\eqref{eq:device-one-body-matrix}, write
$t_{12}=|t_{12}|e^{i\varphi}$. At $\epsilon_1=\epsilon_2=\epsilon_0$
and $|t_{12}|>0$, its eigenstates are
$(\ket{w_1}\pm e^{-i\varphi}\ket{w_2})/\sqrt2$, with energies
$\epsilon_0\mp|t_{12}|$. The equal weights give strong charge
mixing even when $|t_{12}|$ is small compared with gaps to other
states.

This grouping preserves all matrix elements in the selected orbital
space. The Hubbard model in Appendix~\ref{sec:device-hubbard-reduction}
makes additional choices about both the spatial states and the
interaction terms to retain.

\subsection{Solving the many-electron matrix: CI and FCI}
\label{sec:device-ci-fci}

We now solve the eigenvalue problem for $H^{(M)}$ in
Eq.~\eqref{eq:device-projected-hamiltonian}, retaining all its
one-body and two-body matrix elements. At fixed electron number
$N$, choose a basis of Slater determinants $\ket{D_I}$, each
specifying which $N$ spin orbitals are occupied. Expand the
interacting state and solve the matrix eigenvalue problem
\begin{equation}
 \ket{\Psi_{N\alpha}}=\sum_I C_I^{(N\alpha)}\ket{D_I},
 \qquad
 \sum_J\bra{D_I}H^{(M)}\ket{D_J}C_J^{(N\alpha)}
       =E_{N\alpha}C_I^{(N\alpha)}.
 \label{eq:device-ci-eigenproblem}
\end{equation}
Here $\alpha$ labels many-electron eigenstates. This is
\emph{configuration interaction} (CI).

With $M$ spatial orbitals and two spin states for each, there are
\begin{equation}
 \dim\Hilb_N^{(M)}=\binom{2M}{N}
 \label{eq:device-fci-dimension}
\end{equation}
Slater determinants at fixed $N$. \emph{Full configuration
interaction} (FCI) includes them all, or all states in an exactly
invariant symmetry sector containing the state of interest.
Truncated CI uses a selected subset. FCI solves the complete
matrix in the retained orbital space; the remaining errors come
from the orbital truncation, numerical precision, and the
underlying device model~\cite{QTCADManyBody}.

The lowest eigenvalue gives $E_N$. Repeating for different
electron numbers supplies the quantities needed in Section~\ref{sec:electron-number-control}:
\begin{equation}
 N_{\mathrm{eq}}(\mu_{\mathrm R})
   \in\operatorname*{arg\,min}_{N}
          [E_N-\mu_{\mathrm R}N],
 \qquad
 \mu_N=E_N-E_{N-1}.
 \label{eq:device-many-body-charge}
\end{equation}
Here $N_{\mathrm{eq}}$ is the zero-temperature equilibrium charge,
and $\mu_N$ is the energy for adding the $N$th electron.
If successive charge sectors are stable, $\mu_N$ is their
transition threshold; otherwise the minimization may skip a sector.
The set notation allows degeneracy at a transition.
For a localized basis, summing
$\langle c_{i\sigma}^\dagger c_{i\sigma}\rangle$ over the spins
and orbitals assigned to a dot gives its mean electron number.
These mean occupations can be fractional at fixed total $N$.

For a fixed physical Hamiltonian, enlarging a nested orbital
basis can only lower or leave unchanged the FCI ground energy
by the variational principle. Check the convergence of energy
differences such as $J$ and $\mu_N$ separately: monotonic
convergence of each total energy permits either direction of
change in their difference.
